%%%%%%%%%%%%%%%%%%%%%
% 2章
%%%%%%%%%%%%%%%%%%%%%
\chapter{修士論文の書き方} \label{chapter:2}


\section{図と表の例} \label{section:figure_table}

図・表には通し番号と見出し(caption)を付け，本文中で当該の図・表に言及する．また，単位や目盛を正確に記す．
図のタイトルは図の下に，表のタイトルは表の上に書く．

例を図\ref{fig:logo}と表\ref{tab:results}に示す．
第\ref{chapter:2}章の図には図2.1, 図2.2, 図2.3，…のように番号が振られ，
第\ref{chapter:2}章の表には表2.1, 表2.2, 表2.3，…のように，図とは独立に番号が振られる．


\begin{figure}[b]
    \centering
    \includegraphics[scale=0.5]{logo_color.png}
    \caption{情報工学科のロゴ}
    \label{fig:logo}
  \end{figure}


\begin {table}[t]
    \centering
  \caption{表のタイトル}
  \label{tab:results}
  \begin {tabular}{ccc} \hline
     項目1 & 項目2 & 項目3 \\ \hline
    データ1 & データ2 & データ3 \\
    データ1 & データ2 & データ3 \\
    データ1 & データ2 & データ3 \\ \hline
  \end {tabular}
\end {table}


\section{参考文献の書き方}

一例として，和文の著書\cite{suetake}，和文の論文誌\cite{kusano}，英文の編書\cite{fuortes}，
英文の論文誌\cite{rice}，国際会議\cite{guibas}，修士論文\cite{chudai}，電子雑誌\cite{iwama}，Webページ\cite{IPSJ}を，
2ページの参考文献の節に載せる．{\em 参考文献には信頼性が高く，後世に残るものを載せるように注意せよ．}

書くべき情報は以下のとおりである．
\begin{itemize}
\item 和文の著書: 著者，書名，シリーズ名(あれば)，発行所，都市，年．
\item 和文の論文誌: 著者，題名，誌名，巻，号，頁，年．
\item 英文の編書: 編者，書名，発行所，都市，年．
\item 英文の論文誌: 著者，題名，誌名，巻，号，頁，年．
\item 国際会議: 著者，題名，予稿集名，都市，コード等，頁，年．
\item 修士論文: 著者，題名，機関名，年．
\item 電子雑誌: 著者，題名，誌名，巻，号，頁（オンライン），DOI，西暦年．
\item Webページ: 著者，Webページの題名，Webサイトの名称（オンライン）（ただし，著者と同じ場合は省略可），入手先〈URL〉（参照日付）．
\end{itemize}
英語の文献はすべて半角文字で書く．参考文献には本文で引用した文献のみ載せる．
情報処理学会の論文誌の原稿執筆案内\cite{IPSJ}も参考になる．

通し番号は，引用順または著者名のアルファベット順に付ける．
文献の引用のしかたは分野ごとに異なるので，{\em 自己流では書かず，当該分野の論文誌などを参考にする}こと．
